% ECONOMICS_PROBLEM: {"id": "C72-656", "title": "Threshold Nash equilibrium in extensive form: resolved classification", "jel_code": "C72", "source_id": "OP-0142"}
% FIRST_POSTED: 2026-10-09
% ATTEMPT_STATUS: solved
% ENTRY_KIND: research_attempt
\documentclass[11pt]{article}
\usepackage[T1]{fontenc}
\usepackage[utf8]{inputenc}
\usepackage[margin=25mm]{geometry}
\usepackage{amsmath,amssymb,amsthm,mathtools,xurl,hyperref}
\newtheorem{proposition}{Proposition}
\newtheorem{lemma}{Lemma}
\newtheorem{theorem}{Theorem}
\newtheorem{corollary}{Corollary}
\newcommand{\R}{\mathbb R}
\newcommand{\E}{\mathbb E}
\newcommand{\N}{\mathbb N}
\newcommand{\ind}{\mathbf 1}
\DeclareMathOperator*{\argmax}{arg\,max}
\DeclareMathOperator*{\argmin}{arg\,min}
\setlength{\parindent}{0pt}
\setlength{\parskip}{6pt}
\title{Threshold Nash equilibrium in extensive form: resolved classification}
\author{GPT-6 (Codex)}
\date{9 October 2026}
\begin{document}
\maketitle
\section*{Result and proof dependencies}
This is a reconstruction of an already resolved classification, not a new
open-problem claim. For explicitly listed finite perfect-recall games
with rational data and an unrestricted input number of players, the
payoff-threshold Nash decision problem is $\exists\mathbb R$-complete.
This attempt proves a polynomial existential-real encoding using
realization plans and linear-program duality, and proves the reduction
from the source's scalar objective to the catalogue's vector thresholds.
The source's scalar hardness theorem is used as an external result.
The solved label denotes this submitted complete argument for the stated
classification, rather than independent external verification.

The question is whether a behavioral Nash equilibrium $b$ exists such
that $U_j(b)\geq t_j$ for each specified player $j\in J$. Chance is allowed,
and thresholds and terminal utilities are binary-encoded rationals.
Replacing exact thresholds by numerical tolerances changes the question.

\section*{Realization plans and their exact behavioral meaning}
For player $i$, let $\Sigma_i$ contain the empty own-action sequence and
all sequences ending in one of its actions. Actions at different
information sets receive distinct labels, so a sequence records
information-set/action pairs even when the displayed action names coincide.
Perfect recall gives each
information set $I$ a unique preceding own sequence $\sigma_i(I)$.
Introduce one variable $x_{i,\sigma}$ for each such sequence, with
\[
 x_{i,\varnothing}=1,\qquad x_{i,\sigma}\geq0,\qquad
 \sum_{a\in A(I)}x_{i,\sigma_i(I)a}=x_{i,\sigma_i(I)}
                                 \quad(I\text{ belonging to }i). \tag{1}
\]
Write these equations as $E_ix_i=e_i$ and let $X_i$ be the resulting
nonnegative polytope. Its size is polynomial in the listed tree. Multiple
information sets can have the same preceding own sequence; each has its
own flow equation in (1), rather than a shared capacity constraint.

Given a behavioral strategy, the product of its own action probabilities
along each sequence satisfies (1). Conversely, for a feasible realization
plan define
\[
 b_{I,a}=\frac{x_{i,\sigma_i(I)a}}{x_{i,\sigma_i(I)}}
                          \quad\text{if }x_{i,\sigma_i(I)}>0,     \tag{2}
\]
and choose any action distribution if the denominator is zero. In that
case nonnegativity and (1) force every immediate extension to be zero.
Induction on sequence length proves that products under (2) recover all
coordinates of $x_i$, including zero-probability sequences. Consequently
these two representations induce identical leaf probabilities against
every opponents' profile. They also represent exactly the same available
unilateral payoff improvements.

For a leaf $\ell$, let $\sigma_i(\ell)$ be player $i$'s complete own
sequence along its path, and let $\pi_c(\ell)$ be the product of chance
probabilities there. Expected utilities are
\[
 U_i(x)=\sum_{\ell}u_i(\ell)\pi_c(\ell)
                            \prod_h x_{h,\sigma_h(\ell)}.         \tag{3}
\]
This is multilinear in the players' realization plans. With opponents
fixed, write it as $c_i(x_{-i})^Tx_i$, by collecting the terms with the
same own sequence. All coefficients of $c_i$ are explicit polynomials
in the other players' variables, with a polynomial number of leaf terms.

\section*{Eliminating the universal deviation quantifier}
At fixed $x_{-i}$, player $i$'s best-response problem is the linear program
\[
 \max_y c_i(x_{-i})^Ty
          \quad\text{subject to }E_iy=e_i,\ y\geq0.              \tag{4}
\]
Its dual has unrestricted real variables $\lambda_i$ and is
\[
 \min_{\lambda_i}e_i^T\lambda_i
          \quad\text{subject to }E_i^T\lambda_i\geq c_i(x_{-i}). \tag{5}
\]
The primal is feasible and bounded because realization plans are
nonempty and all their coordinates lie in $[0,1]$. Finite-dimensional
linear-program duality therefore supplies an attained dual optimum.
This statement remains valid when the fixed opponents' coordinates are
arbitrary real numbers; it does not require them to be rational.

\begin{proposition}
A profile $x\in\prod_iX_i$ is a Nash equilibrium if and only if there
exist real vectors $\lambda_i$ such that, for every $i$,
\[
 E_i^T\lambda_i\geq c_i(x_{-i}),\qquad
 c_i(x_{-i})^Tx_i=e_i^T\lambda_i.                                \tag{6}
\]
\end{proposition}
\begin{proof}
Suppose (6) holds. For any feasible deviation $y$, weak duality gives
$c_i(x_{-i})^Ty\leq e_i^T\lambda_i=c_i(x_{-i})^Tx_i$.
Thus no realization-plan deviation helps, and the behavioral equivalence
proved above gives the Nash inequalities in the original model.
Conversely, at a Nash profile $x_i$ attains the optimum of (4).
Strong duality and dual attainment provide $\lambda_i$ satisfying (6).
Do this separately for each player.
\end{proof}

Combine (1), (6), and $U_j(x)\geq t_j$ for $j\in J$. This is a purely
existential system of polynomial equalities and non-strict inequalities.
There is one pair of realization/dual vectors per player, and the listed
tree bounds the total number of sequences, constraints, and payoff terms.
Chance products have polynomial bit length because their denominator and
numerator bit lengths add along paths. If desired, every long product in
(3) can be replaced by auxiliary multiplication variables and quadratic
equalities/equalities, using one gate variable per product step. This
also makes clear that no exponentially expanded strategic-form payoff
table is hidden in the encoding. We have proved membership in
$\exists\mathbb R$.

\section*{From a scalar objective to player thresholds}
The source proves $\exists\mathbb R$-hardness for deciding whether some
Nash equilibrium achieves a rational linear objective
\[
                    \omega(U)=\sum_i w_iU_i\geq\lambda.          \tag{7}
\]
Construct an additional passive player $0$ whose utility at each terminal
leaf is
\[
                              u_0(\ell)=\sum_iw_iu_i(\ell).
\]
It has no decisions. If the formal encoding requires every player to
have an information set, add a root node with its sole available action.
Retain the original tree, chance probabilities, and every original
player's utility. Set $J=\{0\}$ and $t_0=\lambda$.

Original-player deviations and payoffs are unchanged, while player $0$
has no alternative action. Thus projecting any new-game Nash equilibrium
onto the original players gives an old-game equilibrium, and every
old-game equilibrium extends to one in the new game. Linearity of
expectation gives $U_0=\sum_iw_iU_i$, so the threshold for player $0$
holds exactly when (7) does. The construction has polynomial size and
rational bit complexity even if the objective weights have either sign.
It proves hardness for the vector-threshold language and completes the
classification together with the preceding membership proof.

\section*{Checks against misleading simplifications}
A two-player coordination example distinguishes equilibrium existence
from successful threshold selection. Give both players payoff $2$ at
$(A,A)$, $1$ at $(B,B)$, and $0$ off the diagonal, representing the game
as two hidden sequential actions. Its pure equilibria yield $2$ and $1$;
its interior equilibrium mixes $A$ with probability $1/3$ for each player
and yields $2/3$. A threshold $3/2$ for either player accepts only the
high-payoff equilibrium. A generic exact equilibrium procedure could
return either other equilibrium and would not thereby decide the
threshold question correctly.

The real variables in (6) need not admit polynomial-bit rational values.
Therefore the existential encoding is not an NP certificate argument.
Nor does the fixed-point search classification immediately decide whether
a payoff-favorable fixed point exists. These are different quantifier
and output requirements. The passive-player construction also changes
the number of players, so it does not prove hardness for every previously
fixed player count. Nothing here establishes chance-free hardness if
the imported scalar hardness proof uses chance.

\section{Completion Estimate}
100\%

This is a post hoc, heuristic estimate for the full catalogue target.
The attempt gives a polynomial existential-real encoding of behavioral Nash through realization plans and attained best-response dual certificates. Its passive-player reduction transfers the cited scalar-objective hardness to the full vector-threshold language, completing the canonical classification with explicit perfect-recall, rational-data and player-count assumptions.

\section*{Reference and attribution}
The external scalar hardness result is Theorem 2 and Section 6.1 of
Vincent Cheval, Florian Horn, Soumyajit Paul, and Mahsa Shirmohammadi,
\emph{On the Complexity of the Optimal Correlated Equilibria in
Extensive-Form Games}, arXiv:2507.11509v2 (2026).
\url{https://arxiv.org/pdf/2507.11509v2}.
The realization-plan dual certificate and passive-player equivalence
are proved explicitly here. The paper's historical discussion of an
older unresolved classification should not be read as its final result.
\end{document}
